\section{Closed immersions and quasi-coherent sheaves}
\label{section-closed-immersions-quasi-coherent}

\noindent
The following lemma finally does for quasi-coherent sheaves on schemes
what Modules, Lemma \ref{modules-lemma-i-star-exact} does for abelian sheaves.
See also the discussion in
Modules, Section \ref{modules-section-closed-immersion}.

\begin{lemma}
\label{lemma-i-star-equivalence}
Let $i : Z \to X$ be a closed immersion of schemes. Let
$\mathcal{I} \subset \mathcal{O}_X$ be the quasi-coherent sheaf of ideals
cutting out $Z$. The functor
$$
i_* :
\QCoh(\mathcal{O}_Z)
\longrightarrow
\QCoh(\mathcal{O}_X)
$$
is exact, fully faithful, with essential image those quasi-coherent
$\mathcal{O}_X$-modules $\mathcal{G}$ such that $\mathcal{I}\mathcal{G} = 0$.
\end{lemma}

\begin{proof}
A closed immersion is quasi-compact and separated, see
Lemmas \ref{lemma-closed-immersion-quasi-compact} and
\ref{lemma-closed-immersion-separated}. Hence
Schemes, Lemma \ref{schemes-lemma-push-forward-quasi-coherent}
applies and the pushforward of a quasi-coherent
sheaf on $Z$ is indeed a quasi-coherent sheaf on $X$.

\medskip\noindent
By Modules, Lemma \ref{modules-lemma-i-star-equivalence}
the functor $i_*$ is fully faithful.

\medskip\noindent
Now we turn to the description of the essential image of the
functor $i_*$. We have $\mathcal{I}(i_*\mathcal{F}) = 0$
for any quasi-coherent $\mathcal{O}_Z$-module, for example by
Modules, Lemma \ref{modules-lemma-i-star-equivalence}.
Next, suppose that $\mathcal{G}$
is any quasi-coherent $\mathcal{O}_X$-module such that
$\mathcal{I}\mathcal{G} = 0$. It suffices to show that the canonical map
$$
\mathcal{G} \longrightarrow i_* i^*\mathcal{G}
$$
is an isomorphism\footnote{This was proved in a more general situation
in the proof of  Modules, Lemma \ref{modules-lemma-i-star-equivalence}.}.
In the case of schemes and quasi-coherent modules, working affine locally
on $X$ and using Lemma \ref{lemma-closed-immersion} and
Schemes, Lemma \ref{schemes-lemma-widetilde-pullback}
it suffices to prove the following algebraic statement: Given a ring
$R$, an ideal $I$ and an $R$-module $N$ such that $IN = 0$ the canonical map
$$
N \longrightarrow N \otimes_R R/I,\quad
n \longmapsto n \otimes 1
$$
is an isomorphism of $R$-modules. Proof of this easy algebra fact is omitted.
\end{proof}

\noindent
Let $i : Z \to X$ be a closed immersion. Because of the lemma above we often,
by abuse of notation, denote $\mathcal{F}$ the sheaf $i_*\mathcal{F}$ on $X$.

\begin{lemma}
\label{lemma-largest-quasi-coherent-subsheaf}
Let $X$ be a scheme. Let $\mathcal{F}$ be a quasi-coherent
$\mathcal{O}_X$-module. Let $\mathcal{G} \subset \mathcal{F}$
be a $\mathcal{O}_X$-submodule. There exists a unique quasi-coherent
$\mathcal{O}_X$-submodule $\mathcal{G}' \subset \mathcal{G}$
with the following property: For every quasi-coherent $\mathcal{O}_X$-module
$\mathcal{H}$ the map
$$
\Hom_{\mathcal{O}_X}(\mathcal{H}, \mathcal{G}')
\longrightarrow
\Hom_{\mathcal{O}_X}(\mathcal{H}, \mathcal{G})
$$
is bijective. In particular $\mathcal{G}'$ is the largest quasi-coherent
$\mathcal{O}_X$-submodule of $\mathcal{F}$ contained in $\mathcal{G}$.
\end{lemma}

\begin{proof}
Let $\mathcal{G}_a$, $a \in A$ be the set of quasi-coherent
$\mathcal{O}_X$-submodules contained in $\mathcal{G}$.
Then the image $\mathcal{G}'$ of
$$
\bigoplus\nolimits_{a \in A} \mathcal{G}_a \longrightarrow \mathcal{F}
$$
is quasi-coherent as the image of a map of quasi-coherent sheaves
on $X$ is quasi-coherent and since a direct sum of quasi-coherent sheaves
is quasi-coherent, see
Schemes, Section \ref{schemes-section-quasi-coherent}.
The module $\mathcal{G}'$ is contained in $\mathcal{G}$. Hence this is the
largest quasi-coherent $\mathcal{O}_X$-module contained in $\mathcal{G}$.

\medskip\noindent
To prove the formula, let $\mathcal{H}$ be a quasi-coherent
$\mathcal{O}_X$-module and let $\alpha : \mathcal{H} \to \mathcal{G}$
be an $\mathcal{O}_X$-module map. The image of the composition
$\mathcal{H} \to \mathcal{G} \to \mathcal{F}$ is quasi-coherent
as the image of a map of quasi-coherent sheaves. Hence it is contained
in $\mathcal{G}'$. Hence $\alpha$ factors through $\mathcal{G}'$
as desired.
\end{proof}

\begin{lemma}
\label{lemma-i-upper-shriek}
Let $i : Z \to X$ be a closed immersion of schemes.
There is a functor\footnote{This is likely nonstandard notation.}
$i^! : \QCoh(\mathcal{O}_X) \to \QCoh(\mathcal{O}_Z)$
which is a right adjoint to $i_*$. (Compare
Modules, Lemma \ref{modules-lemma-i-star-right-adjoint}.)
\end{lemma}

\begin{proof}
Given quasi-coherent $\mathcal{O}_X$-module $\mathcal{G}$ we consider
the subsheaf $\mathcal{H}_Z(\mathcal{G})$ of $\mathcal{G}$ of local sections
annihilated by $\mathcal{I}$. By
Lemma \ref{lemma-largest-quasi-coherent-subsheaf}
there is a canonical largest quasi-coherent $\mathcal{O}_X$-submodule
$\mathcal{H}_Z(\mathcal{G})'$. By construction we have
$$
\Hom_{\mathcal{O}_X}(i_*\mathcal{F}, \mathcal{H}_Z(\mathcal{G})')
=
\Hom_{\mathcal{O}_X}(i_*\mathcal{F}, \mathcal{G})
$$
for any quasi-coherent $\mathcal{O}_Z$-module $\mathcal{F}$.
Hence we can set $i^!\mathcal{G} = i^*(\mathcal{H}_Z(\mathcal{G})')$.
Details omitted.
\end{proof}

\noindent
Using the $1$-to-$1$ corresponding between quasi-coherent sheaves
of ideals and closed subschemes (see
Lemma \ref{lemma-closed-immersion-bijection-ideals})
we can define scheme theoretic intersections and unions
of closed subschemes.

\begin{definition}
\label{definition-scheme-theoretic-intersection-union}
Let $X$ be a scheme. Let $Z, Y \subset X$ be closed subschemes
corresponding to quasi-coherent ideal sheaves
$\mathcal{I}, \mathcal{J} \subset \mathcal{O}_X$.
The {\it scheme theoretic intersection} of $Z$ and $Y$
is the closed subscheme of $X$ cut out by $\mathcal{I} + \mathcal{J}$.
The {\it scheme theoretic union} of $Z$ and $Y$
is the closed subscheme of $X$ cut out by
$\mathcal{I} \cap \mathcal{J}$.
\end{definition}

\begin{lemma}
\label{lemma-scheme-theoretic-intersection}
Let $X$ be a scheme. Let $Z, Y \subset X$ be closed subschemes.
Let $Z \cap Y$ be the scheme theoretic intersection of $Z$ and $Y$.
Then $Z \cap Y \to Z$ and $Z \cap Y \to Y$ are closed immersions
and
$$
\xymatrix{
Z \cap Y \ar[r] \ar[d] & Z \ar[d] \\
Y \ar[r] & X
}
$$
is a cartesian diagram of schemes, i.e., $Z \cap Y = Z \times_X Y$.
\end{lemma}

\begin{proof}
The morphisms $Z \cap Y \to Z$ and $Z \cap Y \to Y$ are closed immersions
by Lemma \ref{lemma-closed-immersion-ideals}.
Let $U = \Spec(A)$ be an affine open of $X$ and let $Z \cap U$ and $Y \cap U$
correspond to the ideals $I \subset A$ and $J \subset A$. Then
$Z \cap Y \cap U$ corresponds to $I + J \subset A$. Since
$A/I \otimes_A A/J = A/(I + J)$ we see that the diagram is
cartesian by our description of fibre products of schemes
in Schemes, Section \ref{schemes-section-fibre-products}.
\end{proof}

\begin{lemma}
\label{lemma-scheme-theoretic-union}
Let $S$ be a scheme. Let $X, Y \subset S$ be closed subschemes.
Let $X \cup Y$ be the scheme theoretic union of $X$ and $Y$.
Let $X \cap Y$ be the scheme theoretic intersection of $X$ and $Y$.
Then $X \to X \cup Y$ and $Y \to X \cup Y$ are closed immersions, there is a
short exact sequence
$$
0 \to \mathcal{O}_{X \cup Y} \to \mathcal{O}_X \times \mathcal{O}_Y
\to \mathcal{O}_{X \cap Y} \to 0
$$
of $\mathcal{O}_S$-modules, and the diagram
$$
\xymatrix{
X \cap Y \ar[r] \ar[d] & X \ar[d] \\
Y \ar[r] & X \cup Y
}
$$
is cocartesian in the category of schemes, i.e.,
$X \cup Y = X \amalg_{X \cap Y} Y$.
\end{lemma}

\begin{proof}
The morphisms $X \to X \cup Y$ and $Y \to X \cup Y$ are closed immersions
by Lemma \ref{lemma-closed-immersion-ideals}. In the short exact sequence
we use the equivalence of Lemma \ref{lemma-i-star-equivalence} to think of
quasi-coherent modules on closed subschemes of $S$ as quasi-coherent modules
on $S$. For the first map in the sequence we use the canonical maps
$\mathcal{O}_{X \cup Y} \to \mathcal{O}_X$ and
$\mathcal{O}_{X \cup Y} \to \mathcal{O}_Y$
and for the second map we use the canonical map
$\mathcal{O}_X \to \mathcal{O}_{X \cap Y}$ and
the negative of the canonical map
$\mathcal{O}_Y \to \mathcal{O}_{X \cap Y}$. Then to check
exactness we may work affine locally.
Let $U = \Spec(A)$ be an affine open of $S$ and let $X \cap U$ and $Y \cap U$
correspond to the ideals $I \subset A$ and $J \subset A$. Then
$(X \cup Y) \cap U$ corresponds to $I \cap J \subset A$
and $X \cap Y \cap U$ corresponds to $I + J \subset A$.
Thus exactness follows from the exactness of
$$
0 \to A/I \cap J \to A/I \times A/J \to A/(I + J) \to 0
$$
To show the diagram is cocartesian, suppose we are given a scheme $T$
and morphisms of schemes $f : X \to T$, $g : Y \to T$ agreeing
as morphisms $X \cap Y \to T$. Goal: Show there exists a unique
morphism $h : X \cup Y \to T$ agreeing with $f$ and $g$.
To construct $h$ we may work affine locally on $X \cup Y$, see
Schemes, Section \ref{schemes-section-glueing-schemes}.
If $s \in X$, $s \not \in Y$, then $X \to X \cup Y$ is
an isomorphism in a neighbourhood of $s$ and it is clear
how to construct $h$. Similarly for $s \in Y$, $s \not \in X$.
For $s \in X \cap Y$ we can pick an affine open
$V = \Spec(B) \subset T$ containing $f(s) = g(s)$.
Then we can choose an affine open $U = \Spec(A) \subset S$
containing $s$ such that $f(X \cap U)$ and $g(Y \cap U)$
are contained in $V$. The morphisms $f|_{X \cap U}$
and $g|_{Y \cap V}$ into $V$ correspond to ring maps
$$
B \to A/I
\quad\text{and}\quad
B \to A/J
$$
which agree as maps into $A/(I + J)$. By the short exact sequence
displayed above there is a unique lift of these ring homomorphism
to a ring map $B \to A/I \cap J$ as desired.
\end{proof}









